% ch46_fibrations.tex — Volume IV Chapter 10
% Retrofitted with full Vol I-III pedagogy: cycle stages, etymology, dialogs.

\chapter{Cartesian and coCartesian Fibrations}

\begin{overviewbox}
The Grothendieck construction in ordinary category theory packages a
pseudofunctor $F : \mathcal{S}^{\mathrm{op}} \to \mathbf{Cat}$ as a single
category $\int F$ together with a \term{fibration} $\int F \to \mathcal{S}$
whose fibres are the categories $F(s)$ and whose ``transport'' is the
functoriality of $F$. The $\infty$-categorical version is
\term{Cartesian} (and \term{coCartesian}) \term{fibrations}: morphisms
$p : \mathcal{X} \to \mathcal{S}$ between quasi-categories satisfying a
combinatorial lifting property that encodes coherent transport up the
fibres.

This chapter introduces the definitions and works through the foundational
examples. The technical heart is the notion of a \term{$p$-Cartesian
morphism}: a $1$-simplex $\alpha : x \to y$ in $\mathcal{X}$ such that
restricting along $\alpha$ gives a homotopy equivalence on mapping
spaces. Cartesian morphisms supply the transport data; Cartesian
fibrations are those $p$ in which enough Cartesian morphisms exist.

Examples organize the chapter:
\begin{itemize}
  \item The \term{family fibration}: for a coCartesian fibration, the
        fibres $\mathcal{X}_s$ form a functor $\mathcal{S} \to
        \mathbf{QCat}$.
  \item The \term{arrow fibration} $\mathcal{C}^{\Delta^1} \to \mathcal{C}$
        (evaluation at the source or target): target evaluation is a
        Cartesian fibration; source evaluation is a coCartesian
        fibration.
  \item The \term{universal Cartesian fibration}: a Cartesian fibration
        $\mathcal{U} \to \mathcal{S}$ classifying functors
        $\mathcal{S}^{\mathrm{op}} \to \mathcal{C}\mathrm{at}_\infty$.
        Existence is the content of Chapter~47's straightening theorem.
\end{itemize}

After this chapter, Chapter~47's straightening/unstraightening equivalence
between Cartesian fibrations and functors $\mathcal{S}^{\mathrm{op}} \to
\mathcal{C}\mathrm{at}_\infty$ closes the loop.
\end{overviewbox}


% ═══════════════════════════════════════════════════════════════════════════
\section{Cartesian Morphisms}
\label{sec:cartesian-morphisms}
% ═══════════════════════════════════════════════════════════════════════════

\subsection{Informally}

Let $p : \mathcal{X} \to \mathcal{S}$ be a map of quasi-categories. A
morphism $\alpha : x \to y$ in $\mathcal{X}$ is \term{$p$-Cartesian} if
it presents $y$ as ``the image of $x$ along $\alpha$ in a universal way
relative to the base $\mathcal{S}$.'' Equivalently: $\alpha$ is
$p$-Cartesian when any morphism $z \to y$ in $\mathcal{X}$ that lies over
a factorization $p(z) \to p(x) \to p(y)$ in $\mathcal{S}$ lifts uniquely
to a factorization $z \to x \to y$ in $\mathcal{X}$.

The intuition is direct from $1$-categorical Grothendieck fibrations:
$\alpha$ is the \term{transport} of $x$ along $p(\alpha)$, and any other
morphism into $y$ should ``factor through'' $\alpha$ in a unique way.

\subsection{Formalizing}

\begin{nameorigin}
\term{Cartesian} (as in ``Cartesian morphism,'' ``Cartesian fibration'')
is from \emph{Grothendieck} via the standard categorical
\term{Cartesian square} (= pullback square). A Cartesian morphism is
one that fits into a universal pullback-style square with respect to
the base; a Cartesian fibration is a family of categories with enough
Cartesian morphisms to encode the contravariant action of the base.
The term ``Cartesian'' goes back to Descartes (1596--1650) via
``Cartesian product'' — the original universal-product construction
in set theory. Grothendieck (1960s) lifted the terminology to fibred
categories; Lurie (HTT 2009) extended it to the $\infty$-categorical
setting.
\end{nameorigin}

\begin{definition}[$p$-Cartesian morphism]
\label{def:p-cartesian}
Let $p : \mathcal{X} \to \mathcal{S}$ be a map of simplicial sets and
$\alpha : x \to y$ a $1$-simplex in $\mathcal{X}$. We say $\alpha$ is
\term{$p$-Cartesian} if for every $n \geq 2$, every commutative diagram
of solid arrows
\begin{equation*}
  \begin{tikzcd}[column sep=large]
    \Lambda^n_n \arrow[r] \arrow[d] & \mathcal{X} \arrow[d, "p"] \\
    \Delta^n \arrow[r] \arrow[ur, dashed] & \mathcal{S}
  \end{tikzcd}
\end{equation*}
in which the restriction of the upper map to the final edge $\Delta^{\{n-1,
n\}}$ of $\Lambda^n_n$ is $\alpha$, admits a (dashed) lift.
\end{definition}

\begin{howtosay}
$\alpha$ \sayit{``a $p$-Cartesian morphism'' or ``$p$-Cartesian''}; the
condition specifies a \emph{lifting} property against outer horns
$\Lambda^n_n \hookrightarrow \Delta^n$ that have $\alpha$ on the final
edge. The dual coCartesian condition uses initial-edge horns
$\Lambda^n_0$.
\end{howtosay}

\begin{theorem}[Equivalent characterizations]
\label{thm:cartesian-equiv}
For $\alpha : x \to y$ in $\mathcal{X}$ over $f : s' \to s$ in
$\mathcal{S}$, the following are equivalent:
\begin{enumerate}
  \item $\alpha$ is $p$-Cartesian.
  \item For every $z \in \mathcal{X}$, the canonical map
        \begin{equation*}
          \mathrm{Map}_\mathcal{X}(z, x) \;\longrightarrow\;
          \mathrm{Map}_\mathcal{X}(z, y) \times_{\mathrm{Map}_\mathcal{S}(p(z), p(y))} \mathrm{Map}_\mathcal{S}(p(z), p(x))
        \end{equation*}
        is a weak homotopy equivalence of Kan complexes.
  \item Every diagram
        \begin{equation*}
          \begin{tikzcd}
            z \arrow[rr, bend left] \arrow[r, dashed, "g"'] & x \arrow[r, "\alpha"] & y
          \end{tikzcd}
        \end{equation*}
        admits a unique (up to coherent homotopy) lift $g$ for any
        compatible $z \to y$ in $\mathcal{X}$ and $p(z) \to p(x)$ in
        $\mathcal{S}$.
\end{enumerate}
\end{theorem}

\begin{remark}[Cartesian for ordinary categories]
For an ordinary functor $p : \mathbf{C} \to \mathbf{S}$ regarded as a map
of nerves, the condition collapses: $\alpha$ is $p$-Cartesian iff for
every $g : z \to y$ in $\mathbf{C}$ and every factorization $p(z) \to p(x)
\to p(y)$ of $p(g)$ in $\mathbf{S}$, there is a unique $\bar g : z \to x$
making both pieces commute. This is the standard $1$-categorical
definition of Grothendieck-Cartesian.
\end{remark}


% ═══════════════════════════════════════════════════════════════════════════
\section{Cartesian Fibrations}
\label{sec:cartesian-fibrations}
% ═══════════════════════════════════════════════════════════════════════════

\subsection{Formalizing}

\begin{nameorigin}
The dichotomy \term{Cartesian} vs.\ \term{coCartesian} reflects the
direction of transport: \emph{Cartesian} fibrations have transport
\emph{against} the base direction (over $f : s \to s'$, transport sends
fibre$_{s'}$ to fibre$_{s}$), so the fibres form a \emph{contravariant}
functor $\mathcal{S}^{\mathrm{op}} \to \mathrm{Cat}_\infty$.
\emph{coCartesian} fibrations have transport \emph{with} the base
direction, so the fibres form a covariant functor. The dual prefix
``co-'' is standard category-theoretic convention.
\end{nameorigin}

\begin{definition}[Cartesian fibration]
\label{def:cartesian-fibration}
A map $p : \mathcal{X} \to \mathcal{S}$ of simplicial sets is a
\term{Cartesian fibration} if:
\begin{itemize}
  \item $p$ is an inner fibration (RLP against inner horns; equivalently,
        the fibres $\mathcal{X}_s = \mathcal{X} \times_\mathcal{S} \{s\}$ are
        quasi-categories).
  \item For every edge $f : s' \to s$ in $\mathcal{S}$ and every $y \in
        \mathcal{X}$ with $p(y) = s$, there exists a $p$-Cartesian
        morphism $\alpha : x \to y$ in $\mathcal{X}$ with $p(\alpha) = f$.
\end{itemize}
\end{definition}

\begin{definition}[coCartesian fibration]
\label{def:cocartesian-fibration}
$p$ is a \term{coCartesian fibration} if it is an inner fibration and
for every $f : s \to s'$ and every $x$ with $p(x) = s$, there exists a
$p$-coCartesian morphism $\alpha : x \to y$ (defined dually) with
$p(\alpha) = f$.
\end{definition}

\begin{example_thm}[Trivial fibrations are Cartesian and coCartesian]
\label{ex:trivial-fib-cart}
A trivial Kan fibration (RLP against all $\partial\Delta^n \hookrightarrow
\Delta^n$) is automatically both Cartesian and coCartesian. The
transport along any base edge is via the lifting property; every morphism
in the total space is both Cartesian and coCartesian.
\end{example_thm}

\begin{example_thm}[The target evaluation $\mathcal{C}^{\Delta^1} \to \mathcal{C}$]
\label{ex:target-eval}
For any quasi-category $\mathcal{C}$, the target evaluation
$\mathrm{ev}_1 : \mathcal{C}^{\Delta^1} \to \mathcal{C}$ (which sends
an arrow $f : a \to b$ to $b$) is a Cartesian fibration. The Cartesian
lift of $g : b' \to b$ over an object $f : a \to b$ is the precomposition
$g^* f = f \circ g : a \to b'$, presented as a $2$-simplex.
\end{example_thm}

\begin{example_thm}[Source evaluation $\mathcal{C}^{\Delta^1} \to \mathcal{C}$]
\label{ex:source-eval}
Dually, $\mathrm{ev}_0 : \mathcal{C}^{\Delta^1} \to \mathcal{C}$ (sending
$f : a \to b$ to $a$) is a \emph{coCartesian} fibration: along $h : a \to
a'$, transport an arrow $f : a \to b$ to its postcomposition $h_* f$.
\end{example_thm}

\begin{nameorigin}
The \term{right fibration} / \term{left fibration} terminology denotes
the directionality of the lifting condition: a right fibration has
\emph{right} lifting against outer horns $\Lambda^n_n$ (final-vertex
horns, on the right), giving the contravariant fibre functor; a left
fibration has right lifting against $\Lambda^n_0$ (initial-vertex
horns), giving the covariant fibre functor. The terminology is Joyal's.
\end{nameorigin}

\begin{example_thm}[Right fibrations]
\label{ex:right-fib}
A \term{right fibration} is a Cartesian fibration in which every morphism
of $\mathcal{X}$ is $p$-Cartesian. Equivalently, $p$ has RLP against
all outer horns $\Lambda^n_n$ — a stronger condition than RLP against
just final-edge horns. Right fibrations are the $\infty$-categorical
analogue of ordinary discrete fibrations (a slice $\mathcal{C}_{/x} \to
\mathcal{C}$ is a right fibration).
\end{example_thm}

\subsection{Reorganizing: fibres are functorial}

\begin{theorem}[Functoriality of fibres]
\label{thm:fibres-functorial}
For a Cartesian fibration $p : \mathcal{X} \to \mathcal{S}$, the assignment
$s \mapsto \mathcal{X}_s$ extends to a functor
\begin{equation*}
  \mathcal{S}^{\mathrm{op}} \;\longrightarrow\; \mathcal{C}\mathrm{at}_\infty
\end{equation*}
(into the $\infty$-category of small $\infty$-categories), called the
\term{functor classified by} $p$. The transport along an edge $f : s' \to
s$ in $\mathcal{S}$ is the functor $f^* : \mathcal{X}_s \to \mathcal{X}_{s'}$
sending $y \mapsto x$, where $\alpha : x \to y$ is the chosen Cartesian
lift.
\end{theorem}

\begin{remark}[Choice issue]
The transport functor $f^*$ depends on choices of Cartesian lifts.
Different choices give homotopic functors, but to make the assignment
\emph{strictly} functorial in $f$ requires the $\infty$-categorical
straightening of Chapter~47.
\end{remark}


% ═══════════════════════════════════════════════════════════════════════════
\section{Examples in the Wild}
\label{sec:fibration-examples}
% ═══════════════════════════════════════════════════════════════════════════

\subsection{The arrow $\infty$-category}

\begin{example_thm}[The arrow $\infty$-category]
\label{ex:arrow-infty-cat}
$\mathcal{C}^{\Delta^1}$ is the $\infty$-category of arrows of
$\mathcal{C}$. Together with $(\mathrm{ev}_0, \mathrm{ev}_1) :
\mathcal{C}^{\Delta^1} \to \mathcal{C} \times \mathcal{C}$, it is the
universal $\infty$-category over $\mathcal{C} \times \mathcal{C}$
classifying ``data of an arrow.''

The product map $(\mathrm{ev}_0, \mathrm{ev}_1)$ is \emph{not} a
Cartesian fibration in general — the joint conditions on source and
target are too restrictive — but each component is a Cartesian
fibration (target) or coCartesian fibration (source).
\end{example_thm}

\subsection{Free fibrations}

\begin{example_thm}[Free fibration generated by a vertex]
\label{ex:free-fib}
For any quasi-category $\mathcal{C}$ and object $x \in \mathcal{C}_0$,
the slice $\mathcal{C}_{/x} \to \mathcal{C}$ is a right fibration. It
classifies the representable presheaf $\mathrm{Map}_\mathcal{C}(-, x) :
\mathcal{C}^{\mathrm{op}} \to \mathcal{S}$ — a presheaf of spaces (Chapter~48).
\end{example_thm}

\subsection{Universal Cartesian fibration}

\begin{remark}[The universal example, foreshadowed]
\label{rem:universal-cart}
There is a Cartesian fibration $\mathcal{U} \to \mathcal{C}\mathrm{at}_\infty^{\mathrm{op}}$
whose fibre over an $\infty$-category $\mathcal{D}$ is $\mathcal{D}$
itself, classifying the identity functor
$\mathrm{id} : \mathcal{C}\mathrm{at}_\infty^{\mathrm{op}} \to
\mathcal{C}\mathrm{at}_\infty$. Every Cartesian fibration $\mathcal{X} \to
\mathcal{S}$ is the pullback of $\mathcal{U}$ along the classifying functor
$\mathcal{S}^{\mathrm{op}} \to \mathcal{C}\mathrm{at}_\infty$. This is
the universal property of $\mathcal{C}\mathrm{at}_\infty$, made precise
by the straightening theorem of Chapter~47.
\end{remark}


% ═══════════════════════════════════════════════════════════════════════════
\section{Exercises}
\label{sec:fibration-exercises}
% ═══════════════════════════════════════════════════════════════════════════

\begin{exercisesbox}
\begin{qa}
\qaitem{Define a $p$-Cartesian morphism.}{A $1$-simplex $\alpha : x \to y$ in $\mathcal{X}$ that lifts every outer horn $\Lambda^n_n$ ending on $\alpha$. Equivalently: restriction along $\alpha$ gives a homotopy equivalence on mapping spaces.}
\qaitem{Why is the lifting condition specifically against $\Lambda^n_n$?}{Because that horn shape leaves the final edge ($\alpha$) free and demands the lift to fit a Cartesian-square diagram of $p(\alpha)$.}
\qaitem{Define a Cartesian fibration.}{An inner fibration $p : \mathcal{X} \to \mathcal{S}$ such that for every edge in $\mathcal{S}$ and every object over its target, there is a $p$-Cartesian morphism over that edge with that target.}
\qaitem{Define a coCartesian fibration.}{Dually: $p$-coCartesian morphisms exist for every edge in $\mathcal{S}$ and object over its source.}
\qaitem{Why are Cartesian fibrations the higher analogue of Grothendieck fibrations?}{Because Cartesian morphisms supply the transport data, and the totality of Cartesian lifts encodes the action of the base on the fibres — functorially up to coherent homotopy.}
\qaitem{What's a right fibration?}{A Cartesian fibration in which every morphism is $p$-Cartesian. Equivalently: RLP against all outer horns $\Lambda^n_n$ (no constraint on which final edge).}
\qaitem{When is $\mathrm{ev}_1 : \mathcal{C}^{\Delta^1} \to \mathcal{C}$ Cartesian?}{Always, for any quasi-category $\mathcal{C}$. The Cartesian lift over $g : b' \to b$ at $f : a \to b$ is the precomposition $f \circ g$.}
\qaitem{When is $\mathrm{ev}_0 : \mathcal{C}^{\Delta^1} \to \mathcal{C}$ coCartesian?}{Always — dually. coCartesian lifts over $h : a \to a'$ at $f : a \to b$ are postcompositions $h^* f$.}
\qaitem{Is $\mathrm{ev}_0$ Cartesian?}{No in general — postcomposing along an arrow in the base is the natural direction; precomposing requires inverses, which may not exist.}
\qaitem{What's a slice as a right fibration?}{$\mathcal{C}_{/x} \to \mathcal{C}$ is a right fibration for any object $x$ of $\mathcal{C}$. It classifies the representable presheaf $\mathrm{Map}_\mathcal{C}(-, x)$.}
\qaitem{What does ``$\mathcal{X}_s$ is the fibre over $s$'' mean?}{$\mathcal{X}_s := \mathcal{X} \times_\mathcal{S} \{s\}$ — the pullback of $\mathcal{X} \to \mathcal{S}$ along $\{s\} \hookrightarrow \mathcal{S}$.}
\qaitem{What's the functor classified by a Cartesian fibration $p$?}{$\mathcal{S}^{\mathrm{op}} \to \mathcal{C}\mathrm{at}_\infty$, $s \mapsto \mathcal{X}_s$. Transport along $f : s' \to s$: $\mathcal{X}_s \to \mathcal{X}_{s'}$ via Cartesian lifts.}
\qaitem{Why is the transport functor only well-defined up to homotopy?}{Because the choice of Cartesian lift over $f$ is only unique up to a contractible space of choices. Strict functoriality requires straightening.}
\qaitem{What is the universal Cartesian fibration?}{$\mathcal{U} \to \mathcal{C}\mathrm{at}_\infty^{\mathrm{op}}$ classifying the identity functor. Every Cartesian fibration is a pullback of $\mathcal{U}$ — content of straightening.}
\qaitem{Difference between Cartesian and inner fibration?}{Inner fibrations have RLP against inner horns only (giving quasi-category fibres but no transport). Cartesian fibrations additionally have the Cartesian lift property (existence of transport data).}
\qaitem{How does Cartesian fibration relate to right fibration?}{Right fibration $\Rightarrow$ Cartesian (in particular). Right fibration: ALL morphisms are Cartesian. General Cartesian: only enough Cartesian lifts to handle base-edge transport.}
\qaitem{What's a left fibration?}{Dual of right: every morphism is $p$-coCartesian. Equivalently, RLP against $\Lambda^n_0$ horns.}
\qaitem{Are right fibrations Kan fibrations?}{Restricted to fibres over a Kan complex, yes — the fibrewise structure forces it. In general, no: right fibrations are RLP against $\Lambda^n_n$ but not $\Lambda^n_k$ for $0 < k < n$.}
\qaitem{What does a Cartesian fibration $p$ over $\Delta^1$ look like?}{Two fibres $\mathcal{X}_0$ and $\mathcal{X}_1$ together with a functor $\mathcal{X}_1 \to \mathcal{X}_0$ given by the Cartesian lifts of the single edge $0 \to 1$.}
\qaitem{And over $\Delta^0$?}{A Cartesian fibration over a point is just a quasi-category $\mathcal{X}_0$ (the single fibre); the Cartesian condition is vacuous.}
\qaitem{Why are Cartesian fibrations the right notion (vs.\ strict pseudofunctors)?}{Because $\infty$-categorical functoriality is strict only up to coherent homotopy; encoding it via Cartesian lifts captures all the coherences without forcing strict commutativity.}
\qaitem{What's the relationship between Cartesian and coCartesian for the same $p$?}{$p^{\mathrm{op}}$ is Cartesian iff $p$ is coCartesian. (Take opposites.) A self-dual map is called a \emph{bifibration} or \emph{biCartesian}.}
\qaitem{What is $\mathcal{C}\mathrm{at}_\infty$?}{The $\infty$-category of small $\infty$-categories. Its objects are quasi-categories; mapping spaces are spaces of functors (equivalences as $1$-cells); the universal Cartesian fibration $\mathcal{U} \to \mathcal{C}\mathrm{at}_\infty^{\mathrm{op}}$ has fibres themselves.}
\end{qa}
\end{exercisesbox}


% ═══════════════════════════════════════════════════════════════════════════
\section*{Chapter Summary}
\addcontentsline{toc}{section}{Chapter Summary}
% ═══════════════════════════════════════════════════════════════════════════

\begin{summarybox}
\textbf{Cartesian morphism.}\quad A $1$-simplex $\alpha$ in $\mathcal{X}$
over $p : \mathcal{X} \to \mathcal{S}$ such that outer horn lifts
$\Lambda^n_n \to \mathcal{X}$ ending on $\alpha$ extend to $\Delta^n$.
Equivalently, restriction along $\alpha$ is a weak equivalence on mapping
spaces.

\textbf{Cartesian and coCartesian fibrations.}\quad Inner fibrations
admitting enough Cartesian (resp.\ coCartesian) lifts. Right fibration:
every morphism Cartesian (so all of $\mathcal{X}$ lies over $\mathcal{S}$
``co-discretely''). Left fibration: every morphism coCartesian.

\textbf{Functoriality of fibres.}\quad A Cartesian fibration $p$ gives
a functor $\mathcal{S}^{\mathrm{op}} \to \mathcal{C}\mathrm{at}_\infty$,
$s \mapsto \mathcal{X}_s$, with transport along edges via Cartesian
lifts. The choice of lifts is unique up to coherent homotopy.

\textbf{Examples.}\quad Target evaluation $\mathcal{C}^{\Delta^1} \to
\mathcal{C}$ (Cartesian); source evaluation (coCartesian); slices
$\mathcal{C}_{/x}$ (right fibration); universal Cartesian fibration
$\mathcal{U} \to \mathcal{C}\mathrm{at}_\infty^{\mathrm{op}}$
(classifying the identity).
\end{summarybox}


% ═══════════════════════════════════════════════════════════════════════════
\section*{Formal Restatement}
\addcontentsline{toc}{section}{Formal Restatement}
% ═══════════════════════════════════════════════════════════════════════════

\begin{formalrestatement}
\textbf{Cartesian morphism.}\quad $\alpha : x \to y$ in $\mathcal{X}$
over $p$ is Cartesian iff $\forall n \geq 2$, $\Lambda^n_n \to \mathcal{X}$
restricting to $\alpha$ on the final edge extends to $\Delta^n$.

\medskip
\textbf{Equivalent characterization.}\quad $\alpha$ is Cartesian iff for
all $z \in \mathcal{X}$:
\begin{equation*}
  \mathrm{Map}_\mathcal{X}(z, x) \to \mathrm{Map}_\mathcal{X}(z, y) \times_{\mathrm{Map}_\mathcal{S}(p(z), p(y))} \mathrm{Map}_\mathcal{S}(p(z), p(x))
\end{equation*}
is a weak equivalence of Kan complexes.

\medskip
\textbf{Cartesian fibration.}\quad $p : \mathcal{X} \to \mathcal{S}$,
inner fibration, with: $\forall (f : s' \to s) \in \mathcal{S}_1, \, \forall y
\in \mathcal{X}_s, \, \exists \alpha \in \mathcal{X}_1 \text{ with } d_0 \alpha = y,
p(\alpha) = f, \alpha \text{ Cartesian}$.

\medskip
\textbf{coCartesian fibration.}\quad Dual: lift over source rather than target.

\medskip
\textbf{Right / left fibration.}\quad Cartesian (resp.\ coCartesian) with
\emph{every} morphism Cartesian (resp.\ coCartesian). Equivalent: RLP
against $\Lambda^n_n$ (resp.\ $\Lambda^n_0$) horns.

\medskip
\textbf{Universal Cartesian fibration.}\quad $\mathcal{U} \to
\mathcal{C}\mathrm{at}_\infty^{\mathrm{op}}$, fibre over $\mathcal{D}$
equals $\mathcal{D}$; classifies $\mathrm{id} :
\mathcal{C}\mathrm{at}_\infty^{\mathrm{op}} \to \mathcal{C}\mathrm{at}_\infty$.
Every Cartesian fibration is a pullback (Chapter~47).
\end{formalrestatement}
